\documentclass[11pt]{article}

\usepackage[margin=1in]{geometry}
\usepackage{amsmath,amssymb,amsthm}
\usepackage{graphicx}
\usepackage[hidelinks]{hyperref}

\newtheorem{theorem}{Theorem}
\newtheorem{lemma}{Lemma}

\title{A Category-Large Class of Norm-Weak USC Points in \(c_0^{**}\)}
\author{fa\_banach\_001, agent\_lane\_17}
\date{2026-06-20}

\begin{document}
\maketitle

\section*{Source Question}

In \emph{Computing sub differential limits of operators on Banach spaces},
arXiv:2212.09650, Rao studies norm-weak upper semicontinuity of the
preduality map. After Proposition 9, the paper asks how large, in the Baire
category sense, the set of norm-weak usc points can be on \(X^{**}\) when
\(X\) is an \(M\)-ideal in its bidual.

\section*{Result}

We give a positive category answer for the prototypical non-reflexive
\(M\)-embedded space \(X=c_0\).

\begin{theorem}
Let \(X=c_0\), so that \(X^{**}=\ell_\infty\) and \(X^{***}=ba(\mathbb N)\).
The set of points of norm-weak upper semicontinuity for the preduality map on
\(X^{**}\) contains a dense open subset of \(\ell_\infty\). In particular it is
residual.
\end{theorem}

\section*{Idea of the Proof}

For \(a\in\ell_\infty\), norm-weak usc is equivalent, by the criterion quoted
in the source paper from Godefroy--Indumathi, to weak-star density of the
predual state space \(S_a\subset \ell_1\) in the full state space
\(S^a\subset ba(\mathbb N)\).

The dense open set is made of sequences whose norm is attained with a strict
gap from the remaining coordinates. For such a sequence, all norming finitely
additive measures live on the Stone--Cech closures of the maximum-coordinate
sets. Atomic measures on the coordinates are weak-star dense in all measures
on those closures. This gives the required state-space density.

\section*{Proof}

We work over the real scalars. Identify \(\ell_\infty\) with \(C(\beta\mathbb
N)\). For \(a=(a_n)\in\ell_\infty\), write \(\widehat a\) for its continuous
extension to \(\beta\mathbb N\), and put \(M=\|a\|_\infty\).

Let \(\mathcal G\) be the set of all nonzero \(a\in\ell_\infty\) such that,
for
\[
E_+=\{n:a_n=M\},\qquad E_-=\{n:a_n=-M\},\qquad E=E_+\cup E_-,
\]
one has \(E\ne\varnothing\) and
\[
  \sup_{n\notin E}|a_n|<M.
\]

\begin{lemma}
Every \(a\in\mathcal G\) is a point of norm-weak usc for the preduality map on
\(\ell_\infty=c_0^{**}\).
\end{lemma}

\begin{proof}
Let \(a\in\mathcal G\). Choose \(\eta>0\) such that \(|a_n|\le M-\eta\) for
all \(n\notin E\). By continuity of \(\widehat a\), the subset of
\(\beta\mathbb N\) on which \(\widehat a=M\) is exactly \(\overline{E_+}^{\beta
\mathbb N}\), and the subset on which \(\widehat a=-M\) is exactly
\(\overline{E_-}^{\beta\mathbb N}\).

The predual state space \(S_a\subset \ell_1\) consists precisely of the signed
atomic measures \(u\) supported on \(E\), positive on \(E_+\), negative on
\(E_-\), and with total variation \(1\). The full state space
\[
S^a=\{\mu\in ba(\mathbb N):\|\mu\|=1,\ \mu(a)=M\}
\]
consists of signed regular measures on \(\beta\mathbb N\) whose positive
variation is supported on \(\overline{E_+}^{\beta\mathbb N}\) and whose negative
variation is supported on \(\overline{E_-}^{\beta\mathbb N}\).

Atomic probability measures on a dense subset of a compact Hausdorff space are
weak-star dense in the probability measures on that compact space. Applying
this separately on \(\overline{E_+}^{\beta\mathbb N}\) and
\(\overline{E_-}^{\beta\mathbb N}\), with the Jordan decomposition of \(\mu\),
shows that every \(\mu\in S^a\) is a weak-star limit of elements of \(S_a\).
Thus \(S_a\) is weak-star dense in \(S^a\), which is equivalent to norm-weak
usc by the source criterion.
\end{proof}

\begin{lemma}
\(\mathcal G\) is open and dense in \(\ell_\infty\).
\end{lemma}

\begin{proof}
Openness follows from the strict gap. If \(a\in\mathcal G\) and
\(\sup_{n\notin E}|a_n|\le M-\eta\), then every sufficiently small
\(\ell_\infty\)-perturbation still has its norm attained on the old maximum
set and remains separated from the remaining coordinates.

For density, fix \(a\in\ell_\infty\) and \(\varepsilon>0\). Let \(M=\|a\|_\infty\).
Choose \(n\) with \(|a_n|>M-\varepsilon/3\). Pick
\(\delta\) with \(M-|a_n|<\delta<\varepsilon\), and define
\[
  b=a+\delta\,\operatorname{sgn}(a_n)e_n,
\]
using either sign if \(a_n=0\). Then \(\|b-a\|_\infty=\delta<\varepsilon\), and
the \(n\)-th coordinate of \(b\) has absolute value \(|a_n|+\delta>M\), while
all other coordinates have absolute value at most \(M\). Hence \(b\in\mathcal
G\).
\end{proof}

The theorem follows immediately from the two lemmas. Since the norm-weak usc
points contain the dense open set \(\mathcal G\), their complement is contained
in the nowhere dense set \(\ell_\infty\setminus\mathcal G\).

\section*{Verification and Limitations}

The proof uses only the source paper's state-space criterion and the standard
identification \(\ell_\infty=C(\beta\mathbb N)\), \(ba(\mathbb N)=C(\beta
\mathbb N)^*\). It does not address arbitrary \(M\)-embedded spaces, nor the
operator-space intermediate question in Remark 13 of the source paper.

\section*{Novelty Check}

Local index search found no exact packet for arXiv:2212.09650. Web searches for
the phrases \texttt{norm-weak usc c0 preduality map},
\texttt{norm-to-weak upper semi-continuity duality c0}, and close variants did
not locate an existing statement of this \(c_0\) category subcase.

\section*{Source Evidence}

\begin{center}
\includegraphics[width=.82\textwidth]{figures/definition_context_page3.png}
\end{center}

\begin{center}
\includegraphics[width=.82\textwidth]{figures/open_problem_context_page9.png}
\end{center}

\section*{References}

\begin{enumerate}
\item T. S. S. R. K. Rao, \emph{Computing sub differential limits of operators
on Banach spaces}, arXiv:2212.09650, 2022.
\item G. Godefroy and V. Indumathi, \emph{Norm-to-weak upper semi-continuity of
the duality and pre-duality mappings}, Set-Valued Analysis 10 (2002), 317--330.
\item P. Harmand, D. Werner, and W. Werner, \emph{\(M\)-ideals in Banach spaces
and Banach algebras}, Lecture Notes in Mathematics 1547, Springer, 1993.
\end{enumerate}

\end{document}
